\section{Discussion}
\label{sec:discussion}

\paragraph{Intelligence is not the bottleneck.}
The load-bearing result is that a bare one-paragraph prompt on the frontier model
already separates rejected from accepted submissions on its own (AUROC
\aurocNaiveFull; \S\ref{sec:res-value}), and the full pipeline does no better. The
scoring parity, read on its face as a negative result for the pipeline, is the
opposite: with or without our prompt the model is already intelligent enough to
produce a trustworthy first-pass reviewing signal---the raw judgment is not what a
deployment is short of, and the reviewer still owns the decision. The open problem is
not building a model that \emph{can} judge, but turning that capability into an
instrument a reviewer can rely on: one
that returns the same verdict on the same paper across runs, grounds its claims in
evidence rather than generating them, and surfaces an anchored, rubric-structured
review instead of a bare number. That is an engineering problem, and it is the one
AIPR addresses. Framed this way the contribution shifts from ``can an LLM judge
quality?''---which our parity result, with a growing literature, answers yes---to
``can a system judge quality \emph{reliably and accountably} at deployment?'',
which is what the reliability gap (Fig.~\ref{fig:naive}b) and the grounded review
speak to. The implication for the field is that effort is better spent on the
processing layer---reliability, grounding, interaction---than on waiting for a more
intelligent base model the task does not appear to need.

\paragraph{The citation dimension: a silent scoring fallback, and what fixes it.}
One dimension does not carry its weight in the graded cohort, and we report the
cause rather than quietly drop it. The citation subscore is pinned at the maximum
on \citationPinnedFull\% of frontier papers (\citationPinnedMini\% on full-mini),
AUROC \citationAUROC{} (chance; \S\ref{sec:res-value}). The mechanism is a scoring
fallback: the grounded citation audit issues literature-search queries and proposes
missing references; when it proposes \emph{none}, the pipeline records a citation
score of 100 (``no gaps found''). In the original frontier grading run the audit
returned an empty recommendation list for every paper---the conservative frontier
auditor proposed nothing---so every paper pinned to 100, and an empty or failed
search is indistinguishable in the score from a genuinely complete bibliography.
We confirmed this by re-grading a \runSDnPapers{}-paper subset on the same pipeline:
the search and audit fire and produce real, varying citation scores, so the pin was
a silent failure mode, not a measurement. The result does not move the validity
finding---citation is the most lightly weighted dimension, and dropping it leaves
the headline unchanged (AUROC \aurocDropCitationFull{} vs.\ \aurocFull{}). It does,
though, point at a general lesson: equipping the model with a live literature-search
tool---letting it look up and \emph{ground} citation recommendations against real
records rather than generate them---is exactly the kind of tool use that turns a
capable model into a more reliable instrument, the same theme as our headline
result. That grounded, search-backed audit is what the deployed product runs; the
fix to the study artifact is to score an empty or failed audit as \emph{unknown}
rather than perfect, and to re-grade under the search-backed configuration.

\paragraph{Relation to concurrent comment-quality benchmarks.}
A concurrent line of work benchmarks the \emph{review comments} humans and AI
systems produce. ReviewBench \citep{reviewbench2026} scores comments by structural
completeness and an LLM-judged ``consequential'' flag, and---like us---concludes
that AI is best framed as complementary to human review. Our study is orthogonal:
ReviewBench measures \emph{intrinsic} properties of review text and, by its
authors' own statement, ``not their correctness,'' whereas we measure the
\emph{extrinsic} validity of a numeric \emph{score} against an external outcome. A
well-structured, claim-targeting comment can still be wrong; a benchmark of comment
form cannot see this, and an outcome-validated score is the complementary evidence.
The two views are additive; a comment-level comparison on the ReviewBench axis,
plus the correctness axis it omits, is left to future work.

\paragraph{Relation to acceptance prediction.}
Predicting editorial outcomes from text is not new: prior work predicts
accept/reject from \emph{review} text and sentiment
\citep{wang2018sentiment,ghosal2019deepsentipeer,kang2018peerread} and analyzes the
drivers of acceptance \citep{jung2025whatdrives}. Our task differs in two ways that
make it both harder and more useful: we score the \emph{manuscript}, not the
reviews (so the signal is available before any review exists), and we validate the
score as a \emph{measurement} against the outcome rather than training a classifier
to maximize accuracy on it. We therefore report agreement and discrimination, not a
tuned acceptance classifier, and we do not claim to predict the acceptance decision.

\paragraph{Contamination and temporal leakage.}
The central threat to any score--outcome study built on a public venue is that the
grading model may have memorized the outcome---the decision or the reviews---during
training, which would inflate the relationship we measure; LLM-generated review
content is already widespread at these venues \citep{liang2024monitoring}. We
control for this by construction in the choice of cohort: the grading model's
knowledge cutoff is August~2025, whereas the primary cohort's (ICLR~2026)
decisions and reviews were released in January~2026, so the outcome could not have
been in training. We add three checks. First, we flag submissions whose arXiv
preprint predates the cutoff and report the headline metrics with and without
them, bounding any residual \emph{paper-text} leakage (the model seeing the
manuscript, as a reviewer would, is not itself a confound; memorizing its
reception is). Second, we run the fully pre-cutoff ICLR~2025 cohort as a
contaminated contrast: if the score--outcome relationship is no stronger there
than on the clean 2026 cohort, memorization is not driving it. Third, a
prospective confirmation grades a future venue's submissions before its decisions
are released, making leakage logically impossible. The arXiv-before-cutoff
sensitivity split and the contaminated contrast are reported in
Appendix~\ref{app:robust}; the prospective confirmation is follow-up. The
distinct question of whether the model's \emph{priors} on authors and
institutions move the score, which matters for the non-blind venues AIPR also
targets, is addressed by a separate prestige-perturbation experiment that holds
the manuscript fixed and varies only the printed identity, reported separately.
We report this threat explicitly; no current LLM-review study resolves it as
directly.

\paragraph{What the human signal's structure implies for our metric.}
Individual reviewers are highly variable, with a bimodal per-reviewer distribution
of claim-undermining comments that persists within accepted and rejected papers
alike \citep{reviewbench2026}. This is why we validate against the \emph{mean}
reviewer rating rather than any single reviewer: a score cannot agree with a ground
truth more reliable than the ground truth is with itself, which is part of why we
lead with the robust low-end claim (H1) rather than with ranking the top. AIPR
also weights novelty and applicability most heavily
(Eq.~\ref{eq:overall})---the contribution-facing axis humans prioritize over the
validity-checking axis easiest for an LLM to produce---and that these dimensions
carry the most signal against the human outcome (Appendix~\ref{app:robust}).
